\section{Alternative Proofs}

This is a collection of alternative proofs for various statements from books, lecture notes etc.

\begin{theorem}[Theorem 10.34 from \cite{lee2011smooth}]
    Let $E$ and $E^\prime$ be smooth vector bundles over a manifold $M$, and let $F \colon E \to E^\prime$ be a bundle homomorphism. Then $\mathrm{ker} F$ and $\mathrm{im} F$ are smooth subbundles of $E$ and $E^\prime$, respectively, if and only if $F$ has constant rank.
\end{theorem}
\begin{proof}
    We present an alternative proof to show that $\mathrm{ker} F$ is a smooth subbundle. Let $\mathrm{rank} F = k$. Choose a local frame $$\sigma_1,\ldots,\sigma_n \colon U \to E,$$ such that $F\circ \sigma_1, \ldots, F \circ \sigma_k$ are linearly independent sections, i.e. a local frame for $\mathrm{im}F$. So for $j > k$ we can write
    \[
    F \circ \sigma_j = \sum_{i = 1}^k \tau_{ij} F \circ \sigma_i
    \]
    for smooth functions \[
    \tau_{ij} \colon U \to \R
    \]
    Define 
    \[
    \sigma_j^\prime := \sigma_j - \sum _{i=1}^k \tau_{ij} \sigma_i
    \]
    Then $F \circ \sigma_j^\prime = 0$, thus $\mathrm{im}\sigma_j^\prime \subseteq \mathrm{ker}F$. Furthermore, $\sigma_{k+1}^\prime,\ldots,\sigma_n^\prime$ are linearly independent, since $\sigma_1,\ldots,\sigma_n$ is a local frame. This shows the statement.
\end{proof}

\begin{proposition}[Proposition 5.17 from \cite{lee2018riemannian}]
    The musical isomorphisms commute with the (total) covariant derivative: if $F$ is any smooth tensor field with a contravariant $i$-th index position and $\flat$ represents the operation of lowering the $i$-th index, then 
    \begin{equation}
        \label{section_alt_proofs:eq:musical_isomorphisms_and_derivative1}
        \nabla_X(F^\flat) = (\nabla_X F)^\flat
    \end{equation}
    for all vector fields $X$. In particular 
    \begin{equation}
        \label{section_alt_proofs:eq:musical_isomorphisms_and_derivative2}
        \nabla (F^\flat) = (\nabla F)^\flat.  
    \end{equation}
    Similiarly, if $G$ has covariant $i$-th position and $\sharp$ denotes raising the $i$-th index, then 
    \begin{equation}
        \label{section_alt_proofs:eq:musical_isomorphisms_and_derivative3}
        \nabla_X(F^\sharp) = (\nabla_X F)^\sharp, \quad  \nabla(F^\sharp) = (\nabla F)^\sharp.
    \end{equation}
\end{proposition}

In the book \cite{lee2018riemannian}, the author uses the identity $F^\flat = \mathrm{tr}(F \otimes g)$, which is however only correct in the case where $F$ is a vector field, since then 
\[
F^\flat = \langle F,- \rangle = \mathrm{tr}(F \otimes g)
\]
by definiton of the trace operator. In general this argument fails, as can be seen by writing each side in a local frame. Thus we present an alternative proof.

\begin{proof}
    First let $F = Z$ be a vector field. Instead of using the identity mentioned above we can also argue as follows: Write $\alpha := \langle Z, - \rangle$. Then for an arbitrary vector field $Y$ we compute:
    \[
    (\nabla_X \alpha) (Y) = X \cdot \alpha(Y) - \alpha(\nabla_X Y) = \langle \nabla_X Z, Y\rangle + \langle Z, \nabla_X Y \rangle - \langle Z, \nabla_X Y \rangle = \langle \nabla_X Z, Y\rangle
    \]
    hence 
    \[
    \nabla_X (Z^\flat) = \nabla_X \langle Z,-\rangle = \langle \nabla_X Z, -\rangle = (\nabla_X Z)^\flat.
    \]
    For general $F$ we may assume (by linearity of the $\flat$ operator)
    \[
    F = Z_1 \otimes \ldots \otimes Z_k \otimes \epsilon^1 \otimes \ldots \otimes \epsilon^l
    \]
    for some vector fields $Z_j$ and covector fields $\epsilon^j$. But then 
    \[
    F^\flat = Z_1 \otimes \ldots \otimes \langle Z_i , -\rangle \otimes \ldots \otimes Z_k \otimes \epsilon^1 \otimes \ldots \otimes \epsilon^l
    \]
    so (\ref{section_alt_proofs:eq:musical_isomorphisms_and_derivative1}) follows from the product rule and what has been shown above. Then (\ref{section_alt_proofs:eq:musical_isomorphisms_and_derivative2}) follows from (\ref{section_alt_proofs:eq:musical_isomorphisms_and_derivative1}) and the definition of the $\flat$ operator as well as the total covariant derivative. For (\ref{section_alt_proofs:eq:musical_isomorphisms_and_derivative3}) we argue as in \cite{lee2018riemannian}: note that $\sharp$ and $\flat$ are inverses to each other (since they operate on the same index position), so by substituting $F = G^\sharp$ in (\ref{section_alt_proofs:eq:musical_isomorphisms_and_derivative1}) and (\ref{section_alt_proofs:eq:musical_isomorphisms_and_derivative2}) and applying $\sharp$ to both sides, the statement follows.
\end{proof}
